\scanpage{60}
% ll. 1–5
On trouve
\[
\begin{aligned}
&a_0+(2+\mu_1)\langle a_0,a'_1\rangle a_1\in D'\qquad\text{i.e.}\\
&\underbrace{\langle a_0,a'_0\rangle}_{\mu_0}+
(2+\mu_1)\underbrace{\langle\deleted{\unsure{a}}a_0,a'_1\rangle
\langle a_1,a'_0\rangle}_{\beta}=0\\
&\text{i.e.}\qquad \mu_0+(2+\mu_1)\beta=0\qquad\text{i.e.}\quad
\deleted{(2+\mu_1)\beta\unsure{=(\text{comme }\beta|(\text{comme }\beta}}\\
&(\text{comme }\mu_0=-\beta)\qquad(1+\mu_1)\beta=0\quad\text{i.e. comme }\beta\ne0,\\
&1+\mu_1=0\quad\text{i.e.}\quad
\deleted{\mu_1=-1\quad\text{i.e.}\quad\lambda_1-1=}\quad
\lambda_1=0,\quad\text{absurde.}
\end{aligned}
\]

% ll. 6–10
Ainsi $\tau_0,\tau_1$ échangent $D_0,D_1$. Prouvons \deleted{\unsure{maintenant}} \inserted{que} $\mu_0=\mu_1=-2$, i.e. $\lambda_0=\lambda_1=-1$, i.e. $\tau_0,\tau_1$ \uline{d'ordre 2} --- i.e. qu'on est dans le cas des polygônes réguliers ordinaires. \deleted{\unsure{Par ce qui précède|Pour ce qui précède}, on aura} \inserted{Prenons coordonnées adaptées :} $D,D'$

\[
\deleted{\unsure{\tau_0\tau_1|\text{Notons}}}\quad
\deleted{D=ka_0,\quad D'=ka_1\quad\text{donc}}
\]

% ll. 11–14
$e_0\in D,\quad e_1=\tau_0e_0$, on aura, si $\tau_0e_1=\zeta e_0$, $\tau_1e_0=\eta e_1$ \inserted{$\tau_1e_1=\xi e_0$}
\[
\tau_0=\begin{pmatrix}0&\zeta\\1&0\end{pmatrix},\quad
\tau_1=\begin{pmatrix}0&\xi\\\eta&0\end{pmatrix}
\]
équations aux valeurs propres
\[
\lambda^2+\zeta=0,\quad\lambda^2+\xi\eta=0
\]
% ll. 15–17
et comme cette équation a la \deleted{\unsure{val. p.}} solution 1, on doit avoir l'autre valeur propre $\lambda_0$ (resp. $\lambda_1$) égale à $-1$\quad O.K.

\begin{center}\rule{1.1cm}{0.3pt}\end{center}

% ll. 18–24
II) Il n'y a pas dans $C=\check{\mathbb P}(V)$ de div. de degré 2 invariant.

Alors on sait que l'image de $G_{0}$ dans $\operatorname{Aut}(C)\simeq \operatorname{GP}(1,k)$ est un des trois groupes $\unsure{\mathcal A|\mathfrak A}_4\subset \unsure{\mathfrak S|\mathscr S}_4,\ \mathcal A_5$. Considérons leurs images inverses dans $\operatorname{SL}(V)$, $\widetilde{\mathcal A}_4\subset\widetilde{\mathfrak S}_4,\ \widetilde{\mathcal A}_5$, \inserted{soit $\widetilde G_0$.} Alors on a donc
\[
G_0\subset\widetilde G_0\cdot\underbrace{k^\times}_{\text{homothéties}}
\]
% archive sheet number: 59
